\chapter{Predictive analyses on IT segment}

This chapter presents obtained modelling results. For each of the models used, in-depth parameter tuning and evaluation 
on both the training-validation and test datasets are performed. 

\section{Evaluation metrics}

Throughout the thesis, particular metrics of goodness of regression are used. Suppose that on the day $t$
we acquired closing price data $Y_t$; therefore, based on the methodology introduced earlier (Subchapter \ref{rolling}),
we obtain a prediction for the day $t+1$ as $\hat{Y}_{t+1}$.

Firstly, in order to measure deviation from the actual $Y_{t+1}$, metric \textbf{absolute percentage error} (APE) is used.
The percentage relativization was chosen in order to unify the scale among all the stocks, making the prediction error comparable
across the stocks and more interpretable. The complete definition of the metric is given as
\begin{equation*}
    \text{APE}(\hat{Y}_{t+1}, Y_{t+1})=\frac{Y_{t+1}-\hat{Y}_{t+1}}{Y_{t+1}}.
\end{equation*}
Building upon the APE definition, the average value of the APE for individual stocks over the entire time interval considered is evaluated.

Apart from that, we introduce a baseline benchmark to measure the model's significance. According to Fama, the asset price developments correspond to 
the random walk \cite{fama1970efficient}. In accordance with this, such naïve method can be introduced, that benchmark predicts the next-day price at the $t+1$
to be equal to the price on the day $t$.

With the benchmark method defined, we quantify performance as the proportion of predictions made by the model that are closer, in terms of APE, to the actual 
closing price than the predictions made by the naïve benchmark. From now on, such benchmarking will be referred as the \textbf{model's dominance}.

\section{Linear regression}

\section{Lasso regularization}

\section{Ridge regularization}

\section{k-Nearest neighbours regression}

\section{Random forest regression}

\section{ARIMA regression}

